\chapter*{Abstract}
Railway networks are complex systems in which infrastructure, train movements, timetables, operational constraints, and 
disruptions interact continuously. In this context, digital twins offer a promising framework for representing, 
simulating, and analysing railway systems. However, constructing a simulation-oriented railway digital twin from 
real-world data remains difficult, mainly because infrastructure and operational datasets are heterogeneous, incomplete, 
and represented at different levels of detail.

This thesis investigates the construction of a railway digital twin prototype for railway network simulation using 
SUMO.\@ The objective is not to build a complete real-time digital twin, but to study how railway infrastructure and 
operational data can be transformed into simulation-ready artefacts. The work proposes a data-driven pipeline that 
processes railway datasets, generates SUMO-compatible networks, reconstructs train trips, executes simulations, and 
analyses the resulting outputs.

Two complementary modelling levels are developed and compared. The macro-level model represents the railway network as a 
station-to-station graph, where stations are modelled as nodes and railway connections as edges. This abstraction enables 
the integration of real punctuality data and supports the reconstruction of train movements on selected parts of the 
network. The micro-level model represents the infrastructure in greater detail through tracks, platforms, switches, 
segmented track sections, and routing structures. This second representation aims to move closer to the physical railway 
infrastructure and to evaluate whether more realistic train simulations can be generated from the available data.

The results show that the proposed pipeline can generate coherent simulation artefacts and that the macro-level model 
provides a robust basis for integrating real operational data into SUMO.\@ The micro-level model successfully reconstructs 
a large part of the infrastructure and supports random-trip simulations, but the integration of real operational data 
remains limited. In particular, the reconstruction of valid paths through the dual graph fails for several station pairs, 
preventing the complete generation of real-data-based micro-level simulations.

This limitation represents one of the main findings of the thesis. It shows that infrastructure coverage alone is not 
sufficient to validate a railway digital twin. A generated network must also be topologically consistent and 
operationally routable. The work therefore highlights the distinction between reconstructing railway elements and 
producing a simulation model capable of supporting real train movements. Overall, this thesis provides a first foundation 
for SUMO-based railway digital twin development and identifies the main challenges that must be addressed to progress 
toward more realistic and operational railway simulations.